\documentclass{../note}

\begin{document}

\begin{zadanie}{1}
Niech $V$ i $W$ będą przestrzeniami liniowymi nad ciałem $K$ i niech $f,g:V\rightarrow W$ będą przekształceniami spełniającymi warunek: dla każdego wektora $v\in V$ istnieje skalar $c_{v}\in K$ taki, że $g(v)=c_{v}f(v)$. Wykaż, że istnieje $c\in K$, dla którego $g=cf$.
\end{zadanie}
\begin{rozwiazanie}
Załóżmy nie wprost, że nie ma jednego takiego $c$, czyli że są dwa niezerowe wektory $v, w \in V$ bez wspólnego $c$, czyli $g(v) = c_vf(v)$, $g(w) = c_wf(w)$ i $c_v \neq c_w$. Zauważmy, że $f(v) \neq 0$ i $f(w) \neq 0$, bo inaczej $0 = c_vf(v) = g(v) = c_wf(v) = 0$ i wtedy jest wspólne $c$. Niech $u = v + w$ i $g(u) = c_uf(u)$. Wtedy
\[c_vf(v) + c_wf(w) = g(v) + g(w) = g(v + w) = g(u) = c_uf(u) = c_uf(v + w) = c_uf(v) + c_uf(w)\]
\[(c_v - c_u)f(v) + (c_w - c_u)f(w) = 0\]
Skoro $c_v \neq c_w$, to jedno z nich nie jest równe $c_u$, czyli z powyższego równania wynika, żę $f(v)$ i $f(w)$ są liniowo zależne, czyli skoro są niezerowe, to $f(v) = af(w)$ dla pewnego niezerowego $a$. Mamy więc:
\[f(v) - af(w) = 0\]
\[f(v - aw) = 0\]
\[c_{v - aw}f(v - aw) = 0\]
\[g(v - aw) = 0\]
\[g(v) - ag(w) = 0\]
\[c_vf(v) = ac_wf(w)\]
\[ac_vf(w) = ac_wf(w)\]
$af(w) \neq 0$, czyli z powyższej równości wynika $c_v = c_w$, sprzeczność.
\end{rozwiazanie}

\begin{zadanie}{2}
Niech $n>1$ oraz $A, B\in M_{n\times n}(\mathbb{C})$. Przypomnijmy, że $r(A)$ oznacza rząd macierzy $A$. Udowodnij lub podaj kontrprzykład:
\begin{enumerate}[label=(\alph*)]
    \item jeśli $A^{2}=B^{2}$, to $A=B$ lub $A=-B$.
    \item jeśli $r(A)=r(B)$, to $r(A^{2})=r(B^{2})$.
    \item jeśli $r(AB)=0$, to $r(BA)=0$.
    \item jeśli $r(AB)=0$, to $r(A)=0$ lub $r(B)=0$.
\end{enumerate}
\end{zadanie}
\begin{rozwiazanie}
\begin{enumerate}[label=(\alph*)]
    \item Niech $A = 
    \begin{pmatrix}
        0 & 1 \\
        0 & 0
    \end{pmatrix}$,
    $B = 
    \begin{pmatrix}
        0 & 0 \\
        1 & 0
    \end{pmatrix}$. Wtedy $A^2 = B^2 = 0$, ale $A \neq B$ i $A \neq -B$.
    \item Niech $A = 
    \begin{pmatrix}
        0 & 1 \\
        0 & 0
    \end{pmatrix}$,
    $B = 
    \begin{pmatrix}
        1 & 0 \\
        1 & 0
    \end{pmatrix}$. Oba są rzędu 1, ale $A^2 = 0$ i $B^2 = B$, czyli $0 = r(A^2) \neq r(B^2) = 1$.
    \item Niech $A = 
    \begin{pmatrix}
        1 & -1 \\
        0 & 0
    \end{pmatrix}$,
    $B = 
    \begin{pmatrix}
        1 & 0 \\
        1 & 0
    \end{pmatrix}$. Wtedy $AB = 0$ i $BA =
    \begin{pmatrix}
        1 & -1 \\
        1 & -1
    \end{pmatrix}$, czyli $0 = r(AB) \neq r(BA) = 1$.
    \item Niech $A = 
    \begin{pmatrix}
        1 & -1 \\
        0 & 0
    \end{pmatrix}$,
    $B = 
    \begin{pmatrix}
        1 & 0 \\
        1 & 0
    \end{pmatrix}$. Wtedy $AB = 0$, czyli $r(AB) = 0$, ale $r(A) = 1$ i $r(B) = 1$.
\end{enumerate}
\end{rozwiazanie}

\end{document}